\documentclass[10pt,a4paper]{amsart}
\usepackage[margin=19mm]{geometry}
\usepackage{amsmath,amssymb,mathtools}
\usepackage[scr=rsfs,cal=euler]{mathalfa}
\usepackage{xfrac}
\usepackage[unicode,hidelinks,bookmarks=false]{hyperref}
\newcommand{\ie}{i.e.\ }
\newcommand{\cf}{cf.\ }
\newcommand{\resp}{respectively\ }
\newcommand{\Subsection}{\subsection}
\providecommand{\nobreakdash}{\nobreak\hbox{-}}
\setlength{\emergencystretch}{2em}
\multlinegap0pt
\usepackage{bm}
\usepackage{bbm}
\usepackage{dsfont}
\usepackage{xspace}
\usepackage{cancel}
\usepackage{mathtools}
\usepackage{amsthm}
\makeatletter
\@ifundefined{theorem}{\newtheorem{theorem}{Theorem}}{}
\makeatother
\title[How to improve the discrimination power of classically simul]{How to improve the discrimination power of classically simulable measurements?}
\author{Yiran Wang, Yongming Li}
\date{Source: arXiv:2607.19070v1 (CC BY 4.0)}
\begin{document}
\maketitle

\noindent\emph{Source and licence.} How to improve the discrimination power of classically simulable measurements?. Version arXiv:2607.19070v1 (CC BY 4.0), distributed under the Creative Commons Attribution 4.0 International licence, as recorded in the arXiv licence metadata for this exact version and shown on the abstract page https://arxiv.org/abs/2607.19070v1. Full rights notice: licenses/arxiv-2607.19070-CC-BY-4.0.txt.\par\smallskip

\noindent\textit{Editorial scope.} Counted unit: the Theorem whose source label is \texttt{theorem 2} in the article (source0.tex lines 402--413, source label \texttt{theorem 2}), delivered together with the article's own complete original proof of it (source0.tex lines 415--492, one \texttt{proof} environment of 78 lines). No mathematics is written, repaired or re-derived: statement and proof are the article's own text line for line. Required context retained uncounted: none. Every definition, display and label of those lines is retained unchanged. Omitted from this package and not redistributed: 1--401, 414--414, 493--1136. Nothing inside a retained excerpt is omitted or altered: every retained body is delivered exactly as the article's lines are written, so no omission receipt of any kind accompanies this package. Citations: the retained excerpts carry no citation command at all, so no bibliography entry of the article is transcribed and no reference is redistributed. Reference adaptation: none was needed and none was made; every reference inside the delivered unit resolves inside it, so no unresolved reference or citation is produced. Printed slips kept literal: no printed word, symbol, hypothesis or number of the counted statement or of its proof was altered. Layout and class: the article's own class is \texttt{revtex4-2} with packages including graphicx, array, amsthm, amsmath, amssymb, mathrsfs, txfonts, pgfplots, algorithm, algorithmic, multirow, mathtools, enumerate, enumitem; the wrapper is the lane's pinned \texttt{amsart} editorial wrapper at 10pt on A4, which restates the theorem-like environments the retained text needs and adds the article's own macro definitions that the retained text needs (none), with extra wrapper packages (bm, bbm, dsfont, xspace, cancel, mathtools). The delivered numbering is the unit's own. Assets: no retained span contains a figure environment, a graphics-inclusion command, a PDF-page inclusion, a table or a TikZ picture; no image file is copied into this package, the package declares no asset dependency and no image file is redistributed with it. Licence: version arXiv:2607.19070v1 of this article is distributed under the Creative Commons Attribution 4.0 International licence, as recorded in the arXiv licence metadata for this exact version and shown on the abstract page https://arxiv.org/abs/2607.19070v1. A full rights notice is delivered at licenses/arxiv-2607.19070-CC-BY-4.0.txt. Characters: every retained excerpt is pure ASCII, so the delivered engine, XeLaTeX with its default Latin Modern text font, reports no missing character.
\begin{theorem}\label{theorem 2}
	Let $\mathcal{M}_{A\rightarrow B}$ be a quantum channel and $\omega$ be a pure magic state. 
	Suppose that there exist a number $\lambda \geqslant 1$ and two CPWP channels $\mathcal{K}$ and $\mathcal{L}$ such that $\mathcal{M} = \lambda \mathcal{K} - (\lambda-1)\mathcal{L}$. 
	Then the exact simulation cost of $\mathcal{M}$ satisfies
	\begin{equation}
		\begin{aligned}
			\frac{M(\mathcal{M})}{M(\omega)} &\leqslant C_{\mathcal{A}}(\mathcal{M};\omega) \\ &\leqslant \min\left\{n:\mathrm{Tr}(E\omega^{\otimes n})=1,0\leqslant W(E|\mathbf{u}) \leqslant \frac{1}{\lambda},  \forall\mathbf{u}\right\},
		\end{aligned}
	\end{equation}
	where $E$ is an effect such that $0 \leqslant E \leqslant \mathds{1}$, and we set the upper bound to $+\infty$ when no such effect $E$ exists.
	
\end{theorem}

\begin{proof}
	To begin with, we prove the lower bound. 
	Suppose that $n$ copies of $\omega$ can exactly simulate $\mathcal{M}$. 
	Then there exists a CPWP channel $\mathcal{N} \in \mathcal{A}$ such that
	\begin{equation}
		\mathcal{M}(\cdot) = \mathcal{N}\left(\omega^{\otimes n}\otimes\cdot\right).
	\end{equation}
	Next, we can get that
	\begin{equation}
			M(\mathcal{M}) \leqslant M (\omega^{\otimes n}) = nM(\omega),
	\end{equation}
	where the inequality follows from the monotonicity of the mana of quantum channels, the equality follows from the additivity of the mana of quantum states.
	Then we can get
	\begin{equation}
		\frac{M(\mathcal{M})}{M(\omega)} \leqslant n.
	\end{equation}
	Due to the definition of the exact simulation cost, we have
	\begin{equation}
		\frac{M(\mathcal{M})}{M(\omega)} \leqslant C_{\mathcal{A}}(\mathcal{M};\omega).
	\end{equation}
	
	We define a linear map $\mathcal{N}_{RA\rightarrow B}$ as follows:
	\begin{equation}\label{definition of N}
		\begin{aligned}
			\mathcal{N}_{RA\rightarrow B}(\rho_{RA}) =& \mathcal{M}_{A\rightarrow B}\Bigl(\mathrm{Tr}_R\left[(E_R\otimes\mathds{1}_A)\rho_{RA}\right]\Bigr) \\ &+ \mathcal{L}_{A\rightarrow B}\Bigl(\mathrm{Tr}_R\left[\bigl((\mathds{1}_R-E_R)\otimes\mathds{1}_A\bigr)\rho_{RA}\right]\Bigr).
		\end{aligned}
	\end{equation}
	Since $E$ is an effect, $\rho$ is a quantum state, it is easy to verify that $\mathcal{N}$ is a completely positive map.
	Next, since $\mathcal{M}$ and $\mathcal{L}$ are trace-preserving maps, we have
	\begin{equation}
		\begin{aligned}
			\mathrm{Tr}\left[\mathcal{N}_{RA\rightarrow B}(\rho_{RA})\right] &= \mathrm{Tr}\left[ (E_R\otimes\mathds{1}_A)\rho_{RA} \right] \\&+ \mathrm{Tr}\left[\bigl((\mathds{1}_R-E_R)\otimes\mathds{1}_A\bigr)\rho_{RA}\right]  \\ &= \mathrm{Tr}(\rho_{RA}).
		\end{aligned}
	\end{equation}
	Thus, $\mathcal{N}$ is a quantum channel.
	
	Subsequently, we prove that $\mathcal{N}$ is a CPWP channel.
	For arbitrary phase space point $\mathbf{u}_R$ and $\mathbf{u}_A$, we have
	\begin{equation}
		\begin{aligned}
			W_{\mathcal{N}}(\mathbf{v}_B|&\mathbf{u}_R,\mathbf{u}_A) = \frac{1}{d_B}\mathrm{Tr}\left[A_{\mathbf{v}_B}\mathcal{N}_{RA\rightarrow B}\left(A_{\mathbf{u}_R}\otimes A_{\mathbf{u}_A}\right)\right] \\ &~~~~~~~~~~~= \frac{1}{d_B}\mathrm{Tr}\left[A_{\mathbf{v}_B}\mathcal{M}_{A\rightarrow B}\left(\mathrm{Tr}_R\left[(E_R\otimes\mathds{1}_A)A_{\mathbf{u}_R}\otimes A_{\mathbf{u}_A}\right]\right)\right] \\ &+ \frac{1}{d_B}\mathrm{Tr}\left[A_{\mathbf{v}_B}\mathcal{L}_{A\rightarrow B}\left(\mathrm{Tr}_R\left[\bigl((\mathds{1}_R-E_R)\otimes\mathds{1}_A\bigr)A_{\mathbf{u}_R}\otimes A_{\mathbf{u}_A}\right]\right)\right] \\ &~~~~~~~~~~~= \frac{1}{d_B}\mathrm{Tr}\left[A_{\mathbf{v}_B}\mathcal{M}_{A\rightarrow B}\left(\mathrm{Tr}(E_RA_{\mathbf{u}_R})A_{\mathbf{u}_A}\right)\right] \\ &~~~~~~~~~~~+ \frac{1}{d_B}\mathrm{Tr}\left[A_{\mathbf{v}_B}\mathcal{L}_{A\rightarrow B}\left([1-\mathrm{Tr}(E_RA_{\mathbf{u}_R})]A_{\mathbf{u}_A}\right)\right] \\ &~~~~~~~~~~~= W(E|\mathbf{u}) W_{\mathcal{M}}(\mathbf{v}_B|\mathbf{u}_A) + [1-W(E|\mathbf{u})]W_{\mathcal{L}}(\mathbf{v}_B|\mathbf{u}_A),
		\end{aligned}
	\end{equation}
	where the first equality and last equality follow from the definition of the Wigner function of quantum channels, the second equality follows from Eq. \eqref{definition of N}.
	Since $\mathcal{M} = \lambda \mathcal{K} - (\lambda-1)\mathcal{L}$, we have 
	\begin{equation}
		\begin{aligned}
			W_{\mathcal{M}}(\mathbf{v}_B|\mathbf{u}_A) &= \frac{1}{d_B}\mathrm{Tr}\left[A_{\mathbf{v}_B}\mathcal{M}(A_{\mathbf{u}_A})\right] \\ &= \frac{1}{d_B}\mathrm{Tr}\left[A_{\mathbf{v}_B}\left(\lambda \mathcal{K}(A_{\mathbf{u}_A}) - (\lambda-1)\mathcal{L} (A_{\mathbf{u}_A})\right)\right] \\ &= \lambda W_{\mathcal{K}}(\mathbf{v}_B|\mathbf{u}_A)- (\lambda-1)W_{\mathcal{L}}(\mathbf{v}_B|\mathbf{u}_A).
		\end{aligned}
	\end{equation}
	Thus, we can get that
	\begin{equation}
		\begin{aligned}
			W_{\mathcal{N}}(\mathbf{v}_B|\mathbf{u}_R,\mathbf{u}_A) &= \lambda W(E|\mathbf{u}) W_{\mathcal{K}}(\mathbf{v}_B|\mathbf{u}_A) \\ &+ [1-\lambda W(E|\mathbf{u})]W_{\mathcal{L}}(\mathbf{v}_B|\mathbf{u}_A).
		\end{aligned}
	\end{equation}
	Therefore, we have $\lambda W(E|\mathbf{u}) \geqslant 0$ and $[1-\lambda W(E|\mathbf{u})] \geqslant 0$ since $0 \leqslant W(E|\mathbf{u}) \leqslant 1/\lambda$.
	Since $\mathcal{K}$ and $\mathcal{L}$ are CPWP channels, their Wigner functions are nonnegative. 
	Hence,
	\begin{equation}
		W_{\mathcal{N}} (\mathbf{v}|\mathbf{u}_R,\mathbf{u}_A) \geqslant 0
	\end{equation}
	for all $\mathbf{u}_R$, $\mathbf{u}_A$, and $\mathbf{v}$. 
	Therefore, $\mathcal{N}$ is a CPWP channel.
	
	At last, since $\mathrm{Tr}\left[(\mathds{1}_R-E_R)\omega_R^{\otimes n}\right]=0$, then for every input state $\rho_A$, we have
	\begin{equation}
		\begin{aligned}
			\mathcal{N}_{RA\rightarrow B}\left(\omega_R^{\otimes n}\otimes\rho_A\right) &= \mathcal{M}_{A\rightarrow B}\left( \mathrm{Tr}_R\left[(E_R\otimes\mathds{1}_A)(\omega_R^{\otimes n}\otimes\rho_A)\right]\right) \\ &= \mathcal{M}_{A\rightarrow B}(\rho_A).
		\end{aligned}
	\end{equation}
	Therefore, $n$ copies of pure magic state $\omega$ can exactly simulate $\mathcal{M}$ by a CPWP channel, which implies
	\begin{equation}
		C_{\mathcal{A}}(\mathcal{M};\omega) \leqslant n.
	\end{equation}
	
	This completes the proof.
\end{proof}

\end{document}
